% To merge into the real deck: replace the next line with
% % !TeX root = raytracing/slides/main.tex
% raytracing_preamble.tex
% Common styling for Ray Tracing presentation

\documentclass[10pt]{beamer}
\usetheme{metropolis}
\usefonttheme{professionalfonts}

% --- Color Definitions ---
\definecolor{PrimaryColor}{RGB}{33,52,72}         % Deep navy
\definecolor{SecondaryColor}{RGB}{84,119,146}     % Desaturated blue-gray
\definecolor{AccentColor}{RGB}{100, 156, 165}       % Soft steel blue

\definecolor{BackgroundColor}{RGB}{245,247,250}   % Very light cool gray/blue-tinted white
\definecolor{TextColor}{RGB}{40,55,70}            % Darker, cool-toned charcoal (slightly less saturated)
\definecolor{LightGray}{RGB}{220,225,230}         % Soft, cool light gray
\definecolor{DarkGray}{RGB}{70,90,105}            % Cool-toned dark gray (coherent with Primary/Secondary)

\definecolor{RayColor}{RGB}{230,150,80}           % Muted warm orange (less saturated than pure orange)
\definecolor{ObjectColor}{RGB}{128,101,160}       % Muted violet (adjusted purple to match cool palette)
\definecolor{LightColor}{RGB}{235,200,100}        % Soft golden yellow (to complement cool blues)

% --- Theme Customization ---
\setbeamercolor{background canvas}{bg=BackgroundColor}
\setbeamercolor{normal text}{fg=TextColor}
\setbeamercolor{frametitle}{bg=PrimaryColor, fg=white}
\setbeamercolor{section in toc}{fg=PrimaryColor}
\setbeamercolor{block title}{bg=PrimaryColor!80, fg=white}
\setbeamercolor{block body}{bg=PrimaryColor!10}
\setbeamercolor{alerted text}{fg=AccentColor}
\setbeamercolor{itemize item}{fg=PrimaryColor}
\setbeamercolor{itemize subitem}{fg=SecondaryColor}
\setbeamerfont{frametitle}{size=\large,series=\bfseries}

% --- Packages ---
\usepackage[utf8]{inputenc}
\usepackage{url}
\usepackage{booktabs}
\usepackage{amsmath, amssymb}
\usepackage{fontawesome5}
\usepackage{pifont}
\usepackage[most]{tcolorbox}
\tcbuselibrary{skins}
\usepackage{colortbl}
\usepackage{array}
\usepackage{tikz}
\usepackage{graphicx}
\usetikzlibrary{shapes.callouts, positioning, arrows.meta, shapes.geometric, shadows, calc, patterns, 3d, backgrounds, shadings}
\usepackage{adjustbox}
\usepackage{ragged2e}
\usepackage{pgfplots}
\usepackage{graphicx}
\usepackage{caption}
\usepackage{minted}
\usepackage{subcaption}
\pgfplotsset{compat=1.18}

% --- Custom Commands ---
\newcommand{\cmark}{\textcolor{SecondaryColor}{\ding{51}}}
\newcommand{\xmark}{\textcolor{AccentColor}{\ding{55}}}
\newcommand{\highlight}[1]{\textcolor{PrimaryColor}{\textbf{#1}}}
\newcommand{\raycolor}[1]{\textcolor{RayColor}{\textbf{#1}}}
\newcommand{\objectcolor}[1]{\textcolor{ObjectColor}{\textbf{#1}}}

% --- TikZ styles for ray tracing visualizations ---
\tikzstyle{process} = [rectangle, rounded corners=3mm, minimum width=2cm, minimum height=0.8cm, text centered, draw=PrimaryColor, thick, fill=PrimaryColor!15, drop shadow]
\tikzstyle{arrow} = [thick, PrimaryColor, ->, >=stealth]
\tikzstyle{ray} = [thick, RayColor, ->, >=stealth]
\tikzstyle{lightray} = [thick, LightColor, ->, >=stealth]
\tikzstyle{reflectray} = [thick, SecondaryColor, ->, >=stealth]
\tikzstyle{refractray} = [thick, AccentColor, ->, >=stealth]
\tikzstyle{shadowray} = [thick, DarkGray, ->, >=stealth, dashed]

% --- 3D object styles ---
\tikzstyle{sphere} = [circle, minimum size=1.5cm, draw=ObjectColor, thick, fill=ObjectColor!20, drop shadow]
\tikzstyle{plane} = [rectangle, minimum width=3cm, minimum height=0.2cm, draw=ObjectColor, thick, fill=ObjectColor!20]
\tikzstyle{triangle} = [regular polygon, regular polygon sides=3, minimum size=1.5cm, draw=ObjectColor, thick, fill=ObjectColor!20]

% --- Eye/Camera styles ---
\tikzstyle{eye} = [circle, minimum size=0.8cm, draw=PrimaryColor, thick, fill=PrimaryColor!30]
\tikzstyle{pixel} = [rectangle, minimum size=0.2cm, draw=AccentColor, fill=AccentColor!30]

% --- Text box styles ---
\tikzstyle{conceptbox} = [rectangle, rounded corners, fill=PrimaryColor!10, draw=PrimaryColor, thick, text width=0.8\textwidth, inner sep=8pt]
\tikzstyle{formulabox} = [rectangle, rounded corners, fill=SecondaryColor!10, draw=SecondaryColor, thick, inner sep=6pt]

% --- Custom environments ---
\newtcolorbox{raybox}[1]{
  colback=RayColor!10,
  colframe=RayColor,
  title=#1,
  fonttitle=\bfseries,
  sharp corners
}

\newtcolorbox{conceptbox}[1]{
  colback=PrimaryColor!10,
  colframe=PrimaryColor,
  title=#1,
  fonttitle=\bfseries,
  rounded corners
}

\newtcolorbox{mathbox}[1]{
  colback=SecondaryColor!10,
  colframe=SecondaryColor,
  title=#1,
  fonttitle=\bfseries,
  rounded corners
}

\newenvironment{timeline}{%
  \begin{tikzpicture}[scale=0.8]
    \coordinate (start) at (0,0);
    \newcounter{timelineitem}
    \setcounter{timelineitem}{0}
  }{%
  \end{tikzpicture}
}

\newcommand{\timelineitem}[4]{%
  \stepcounter{timelineitem}%
  % Name of this item’s node
  \edef\thisID{time\thetimelineitem}%
  \ifnum\value{timelineitem}=1
  % First item: anchor at (start)
  \node[rectangle, draw, fill=PrimaryColor!20,
  text width=1.5cm, minimum height=0.8cm, align=center]
  (\thisID) at (start) {\textbf{#2}};
  \else
  % Later items: below=#1 of previous time node
  \pgfmathtruncatemacro\prev{\value{timelineitem}-1}%
  \node[rectangle, draw, fill=PrimaryColor!20,
    text width=1.5cm, minimum height=0.8cm, align=center,
  below=#1 of time\prev]
  (\thisID) {\textbf{#2}};
  \fi
  % Description always to the right of this time node
  \node[rectangle, draw, fill=SecondaryColor!10,
    text width=4cm, minimum height=0.8cm,
  align=left, right=0.3cm of \thisID]
  (desc\thetimelineitem) {\textbf{#3}\\#4};
  % Draw arrow back to previous if not the first
  \ifnum\value{timelineitem}>1
  \draw[->, thick, PrimaryColor]
  (time\prev.south) -- (\thisID.north);
  \fi
}

\tikzset{
  camera/.style={fill=PrimaryColor!60, draw=PrimaryColor!80, rectangle, minimum size=8pt},
  image plane/.style={fill=AccentColor!10, draw=AccentColor!50, opacity=0.8},
  pixel/.style={fill=AccentColor!60, thick},
  primary ray/.style={->, very thick, red!90},
  object/.style={fill=ObjectColor!60, draw=ObjectColor!80, circle, minimum size=12pt},
  fovangle/.style={<->, thick, PrimaryColor, dashed}
}

\tikzset{
  lens/.style={thick, PrimaryColor, line width=3pt},
  focal plane/.style={thick, AccentColor},
  object ray/.style={->, thick, ObjectColor},
  image ray/.style={->, thick, SecondaryColor},
  optical axis/.style={dashed, gray}
}

\newcommand{\lens}[3]{%
  \begingroup
  % half‐dimensions
  \pgfmathsetlengthmacro{\a}{#2}%
  \pgfmathsetlengthmacro{\b}{#3}%
  \pgfmathsetlengthmacro{\c}{#3/2}%
  \begin{scope}[shift={#1}]
    \draw[line join=round, fill=blue!15]
    (0,-{\c})
    arc(-30:30:{\a} and {\b})
    arc(150:210:{\a} and {\b})
    ;
  \end{scope}
  \endgroup
}

% ================================================================
% Reference theme, rebuilt to match combined/deck.pdf
% Rust title bar, cream ground, Adventor headings, thin footer.
% Placed last so it overrides metropolis.
% ================================================================
\usepackage{tgheros}
\usepackage{tgadventor}
\renewcommand{\sfdefault}{qhv}                       % Heros for body copy
\newcommand{\headfont}{\fontfamily{qag}\selectfont}  % Adventor for headings

\definecolor{ThemeRust}{HTML}{9B4632}
\definecolor{ThemeCream}{HTML}{F8F6F2}
\definecolor{ThemeSage}{HTML}{C6CCB1}
\definecolor{ThemeStone}{HTML}{E7E7DB}
\definecolor{ThemeMuted}{HTML}{8B8A88}
\definecolor{ThemeInk}{HTML}{191919}

\setbeamercolor{background canvas}{bg=ThemeCream}
\setbeamercolor{normal text}{fg=ThemeInk}
\setbeamercolor{frametitle}{bg=ThemeRust, fg=ThemeCream}
\setbeamercolor{footer}{bg=ThemeCream, fg=ThemeMuted}
\setbeamercolor{alerted text}{fg=ThemeRust}
\setbeamercolor{itemize item}{fg=ThemeRust}
\setbeamercolor{itemize subitem}{fg=ThemeRust}
\setbeamercolor{palette primary}{bg=ThemeRust, fg=ThemeCream}  % section pages

% full-bleed rust bar across the top
\setbeamertemplate{frametitle}{%
  \nointerlineskip%
  \begin{beamercolorbox}[wd=\paperwidth,sep=0pt]{frametitle}%
    \vbox{%
      \vskip0.60em
      \hbox{\hskip0.9em%
        \begin{minipage}{\dimexpr\paperwidth-1.8em\relax}%
          \raggedright\headfont\bfseries\large\insertframetitle%
        \end{minipage}}%
      \vskip0.55em
    }%
  \end{beamercolorbox}%
}

% thin three-zone footer
\setbeamertemplate{footline}{%
  \hbox{%
  \begin{beamercolorbox}[wd=.40\paperwidth,ht=1.9ex,dp=1.0ex,left]{footer}%
    \hspace*{0.9em}\tiny\insertshortauthor
  \end{beamercolorbox}%
  \begin{beamercolorbox}[wd=.28\paperwidth,ht=1.9ex,dp=1.0ex,center]{footer}%
    \tiny\insertshorttitle
  \end{beamercolorbox}%
  \begin{beamercolorbox}[wd=.32\paperwidth,ht=1.9ex,dp=1.0ex,right]{footer}%
    \tiny\insertframenumber\,/\,\inserttotalframenumber\hspace*{0.9em}
  \end{beamercolorbox}}%
  \vskip1pt
}

% full-bleed rust title page
\setbeamertemplate{title page}{%
  \begin{tikzpicture}[remember picture, overlay]
    \fill[ThemeRust] (current page.south west) rectangle (current page.north east);
  \end{tikzpicture}%
  \vfill
  \begin{flushright}
    {\headfont\bfseries\LARGE\color{ThemeCream}\inserttitle\par}
    \vspace{0.3em}
    {\color{ThemeCream!85}\small\insertsubtitle\par}
    \vspace{1.3em}
    {\color{ThemeCream}\small\insertauthor\par}
    \vspace{0.5em}
    {\color{ThemeCream!80}\scriptsize\insertinstitute\par}
    \vspace{0.8em}
    {\color{ThemeCream!75}\scriptsize\insertdate\par}
  \end{flushright}
  \vfill
}

% --- Deck palette remapped onto the theme -------------------------------
% Overrides the navy/violet set defined at the top of this file. Every TikZ
% figure and tcolorbox references these names, so they all follow along.
\definecolor{PrimaryColor}{HTML}{9B4632}     % rust: highlights, frames, titles
\definecolor{SecondaryColor}{HTML}{7A5C4B}   % warm brown: secondary boxes
\definecolor{AccentColor}{HTML}{7E8C57}      % sage: arrows, alerted text
\definecolor{ObjectColor}{HTML}{A9745B}      % clay: solid shapes
\definecolor{LightColor}{HTML}{D9A441}       % warm gold
\definecolor{TextColor}{HTML}{191919}
\definecolor{BackgroundColor}{HTML}{F8F6F2}
\definecolor{LightGray}{HTML}{E7E7DB}
\definecolor{DarkGray}{HTML}{6B6257}

% the four fern rules -- must match COLS in part 3/figs5.py
\definecolor{RuleA}{HTML}{C0392B}
\definecolor{RuleB}{HTML}{2E7D32}
\definecolor{RuleC}{HTML}{B8860B}
\definecolor{RuleD}{HTML}{1F6F8B}

  (the teacher's actual preamble)
% !TeX root = raytracing/slides/main.tex
% raytracing_preamble.tex
% Common styling for Ray Tracing presentation

\documentclass[10pt]{beamer}
\usetheme{metropolis}
\usefonttheme{professionalfonts}

% --- Color Definitions ---
\definecolor{PrimaryColor}{RGB}{33,52,72}         % Deep navy
\definecolor{SecondaryColor}{RGB}{84,119,146}     % Desaturated blue-gray
\definecolor{AccentColor}{RGB}{100, 156, 165}       % Soft steel blue

\definecolor{BackgroundColor}{RGB}{245,247,250}   % Very light cool gray/blue-tinted white
\definecolor{TextColor}{RGB}{40,55,70}            % Darker, cool-toned charcoal (slightly less saturated)
\definecolor{LightGray}{RGB}{220,225,230}         % Soft, cool light gray
\definecolor{DarkGray}{RGB}{70,90,105}            % Cool-toned dark gray (coherent with Primary/Secondary)

\definecolor{RayColor}{RGB}{230,150,80}           % Muted warm orange (less saturated than pure orange)
\definecolor{ObjectColor}{RGB}{128,101,160}       % Muted violet (adjusted purple to match cool palette)
\definecolor{LightColor}{RGB}{235,200,100}        % Soft golden yellow (to complement cool blues)

% --- Theme Customization ---
\setbeamercolor{background canvas}{bg=BackgroundColor}
\setbeamercolor{normal text}{fg=TextColor}
\setbeamercolor{frametitle}{bg=PrimaryColor, fg=white}
\setbeamercolor{section in toc}{fg=PrimaryColor}
\setbeamercolor{block title}{bg=PrimaryColor!80, fg=white}
\setbeamercolor{block body}{bg=PrimaryColor!10}
\setbeamercolor{alerted text}{fg=AccentColor}
\setbeamercolor{itemize item}{fg=PrimaryColor}
\setbeamercolor{itemize subitem}{fg=SecondaryColor}
\setbeamerfont{frametitle}{size=\large,series=\bfseries}

% --- Packages ---
\usepackage[utf8]{inputenc}
\usepackage{url}
\usepackage{booktabs}
\usepackage{amsmath, amssymb}
\usepackage{fontawesome5}
\usepackage{pifont}
\usepackage[most]{tcolorbox}
\tcbuselibrary{skins}
\usepackage{colortbl}
\usepackage{array}
\usepackage{tikz}
\usepackage{graphicx}
\usetikzlibrary{shapes.callouts, positioning, arrows.meta, shapes.geometric, shadows, calc, patterns, 3d, backgrounds, shadings}
\usepackage{adjustbox}
\usepackage{ragged2e}
\usepackage{pgfplots}
\usepackage{graphicx}
\usepackage{caption}
\usepackage{minted}
\usepackage{subcaption}
\pgfplotsset{compat=1.18}

% --- Custom Commands ---
\newcommand{\cmark}{\textcolor{SecondaryColor}{\ding{51}}}
\newcommand{\xmark}{\textcolor{AccentColor}{\ding{55}}}
\newcommand{\highlight}[1]{\textcolor{PrimaryColor}{\textbf{#1}}}
\newcommand{\raycolor}[1]{\textcolor{RayColor}{\textbf{#1}}}
\newcommand{\objectcolor}[1]{\textcolor{ObjectColor}{\textbf{#1}}}

% --- TikZ styles for ray tracing visualizations ---
\tikzstyle{process} = [rectangle, rounded corners=3mm, minimum width=2cm, minimum height=0.8cm, text centered, draw=PrimaryColor, thick, fill=PrimaryColor!15, drop shadow]
\tikzstyle{arrow} = [thick, PrimaryColor, ->, >=stealth]
\tikzstyle{ray} = [thick, RayColor, ->, >=stealth]
\tikzstyle{lightray} = [thick, LightColor, ->, >=stealth]
\tikzstyle{reflectray} = [thick, SecondaryColor, ->, >=stealth]
\tikzstyle{refractray} = [thick, AccentColor, ->, >=stealth]
\tikzstyle{shadowray} = [thick, DarkGray, ->, >=stealth, dashed]

% --- 3D object styles ---
\tikzstyle{sphere} = [circle, minimum size=1.5cm, draw=ObjectColor, thick, fill=ObjectColor!20, drop shadow]
\tikzstyle{plane} = [rectangle, minimum width=3cm, minimum height=0.2cm, draw=ObjectColor, thick, fill=ObjectColor!20]
\tikzstyle{triangle} = [regular polygon, regular polygon sides=3, minimum size=1.5cm, draw=ObjectColor, thick, fill=ObjectColor!20]

% --- Eye/Camera styles ---
\tikzstyle{eye} = [circle, minimum size=0.8cm, draw=PrimaryColor, thick, fill=PrimaryColor!30]
\tikzstyle{pixel} = [rectangle, minimum size=0.2cm, draw=AccentColor, fill=AccentColor!30]

% --- Text box styles ---
\tikzstyle{conceptbox} = [rectangle, rounded corners, fill=PrimaryColor!10, draw=PrimaryColor, thick, text width=0.8\textwidth, inner sep=8pt]
\tikzstyle{formulabox} = [rectangle, rounded corners, fill=SecondaryColor!10, draw=SecondaryColor, thick, inner sep=6pt]

% --- Custom environments ---
\newtcolorbox{raybox}[1]{
  colback=RayColor!10,
  colframe=RayColor,
  title=#1,
  fonttitle=\bfseries,
  sharp corners
}

\newtcolorbox{conceptbox}[1]{
  colback=PrimaryColor!10,
  colframe=PrimaryColor,
  title=#1,
  fonttitle=\bfseries,
  rounded corners
}

\newtcolorbox{mathbox}[1]{
  colback=SecondaryColor!10,
  colframe=SecondaryColor,
  title=#1,
  fonttitle=\bfseries,
  rounded corners
}

\newenvironment{timeline}{%
  \begin{tikzpicture}[scale=0.8]
    \coordinate (start) at (0,0);
    \newcounter{timelineitem}
    \setcounter{timelineitem}{0}
  }{%
  \end{tikzpicture}
}

\newcommand{\timelineitem}[4]{%
  \stepcounter{timelineitem}%
  % Name of this item’s node
  \edef\thisID{time\thetimelineitem}%
  \ifnum\value{timelineitem}=1
  % First item: anchor at (start)
  \node[rectangle, draw, fill=PrimaryColor!20,
  text width=1.5cm, minimum height=0.8cm, align=center]
  (\thisID) at (start) {\textbf{#2}};
  \else
  % Later items: below=#1 of previous time node
  \pgfmathtruncatemacro\prev{\value{timelineitem}-1}%
  \node[rectangle, draw, fill=PrimaryColor!20,
    text width=1.5cm, minimum height=0.8cm, align=center,
  below=#1 of time\prev]
  (\thisID) {\textbf{#2}};
  \fi
  % Description always to the right of this time node
  \node[rectangle, draw, fill=SecondaryColor!10,
    text width=4cm, minimum height=0.8cm,
  align=left, right=0.3cm of \thisID]
  (desc\thetimelineitem) {\textbf{#3}\\#4};
  % Draw arrow back to previous if not the first
  \ifnum\value{timelineitem}>1
  \draw[->, thick, PrimaryColor]
  (time\prev.south) -- (\thisID.north);
  \fi
}

\tikzset{
  camera/.style={fill=PrimaryColor!60, draw=PrimaryColor!80, rectangle, minimum size=8pt},
  image plane/.style={fill=AccentColor!10, draw=AccentColor!50, opacity=0.8},
  pixel/.style={fill=AccentColor!60, thick},
  primary ray/.style={->, very thick, red!90},
  object/.style={fill=ObjectColor!60, draw=ObjectColor!80, circle, minimum size=12pt},
  fovangle/.style={<->, thick, PrimaryColor, dashed}
}

\tikzset{
  lens/.style={thick, PrimaryColor, line width=3pt},
  focal plane/.style={thick, AccentColor},
  object ray/.style={->, thick, ObjectColor},
  image ray/.style={->, thick, SecondaryColor},
  optical axis/.style={dashed, gray}
}

\newcommand{\lens}[3]{%
  \begingroup
  % half‐dimensions
  \pgfmathsetlengthmacro{\a}{#2}%
  \pgfmathsetlengthmacro{\b}{#3}%
  \pgfmathsetlengthmacro{\c}{#3/2}%
  \begin{scope}[shift={#1}]
    \draw[line join=round, fill=blue!15]
    (0,-{\c})
    arc(-30:30:{\a} and {\b})
    arc(150:210:{\a} and {\b})
    ;
  \end{scope}
  \endgroup
}

% ================================================================
% Reference theme, rebuilt to match combined/deck.pdf
% Rust title bar, cream ground, Adventor headings, thin footer.
% Placed last so it overrides metropolis.
% ================================================================
\usepackage{tgheros}
\usepackage{tgadventor}
\renewcommand{\sfdefault}{qhv}                       % Heros for body copy
\newcommand{\headfont}{\fontfamily{qag}\selectfont}  % Adventor for headings

\definecolor{ThemeRust}{HTML}{9B4632}
\definecolor{ThemeCream}{HTML}{F8F6F2}
\definecolor{ThemeSage}{HTML}{C6CCB1}
\definecolor{ThemeStone}{HTML}{E7E7DB}
\definecolor{ThemeMuted}{HTML}{8B8A88}
\definecolor{ThemeInk}{HTML}{191919}

\setbeamercolor{background canvas}{bg=ThemeCream}
\setbeamercolor{normal text}{fg=ThemeInk}
\setbeamercolor{frametitle}{bg=ThemeRust, fg=ThemeCream}
\setbeamercolor{footer}{bg=ThemeCream, fg=ThemeMuted}
\setbeamercolor{alerted text}{fg=ThemeRust}
\setbeamercolor{itemize item}{fg=ThemeRust}
\setbeamercolor{itemize subitem}{fg=ThemeRust}
\setbeamercolor{palette primary}{bg=ThemeRust, fg=ThemeCream}  % section pages

% full-bleed rust bar across the top
\setbeamertemplate{frametitle}{%
  \nointerlineskip%
  \begin{beamercolorbox}[wd=\paperwidth,sep=0pt]{frametitle}%
    \vbox{%
      \vskip0.60em
      \hbox{\hskip0.9em%
        \begin{minipage}{\dimexpr\paperwidth-1.8em\relax}%
          \raggedright\headfont\bfseries\large\insertframetitle%
        \end{minipage}}%
      \vskip0.55em
    }%
  \end{beamercolorbox}%
}

% thin three-zone footer
\setbeamertemplate{footline}{%
  \hbox{%
  \begin{beamercolorbox}[wd=.40\paperwidth,ht=1.9ex,dp=1.0ex,left]{footer}%
    \hspace*{0.9em}\tiny\insertshortauthor
  \end{beamercolorbox}%
  \begin{beamercolorbox}[wd=.28\paperwidth,ht=1.9ex,dp=1.0ex,center]{footer}%
    \tiny\insertshorttitle
  \end{beamercolorbox}%
  \begin{beamercolorbox}[wd=.32\paperwidth,ht=1.9ex,dp=1.0ex,right]{footer}%
    \tiny\insertframenumber\,/\,\inserttotalframenumber\hspace*{0.9em}
  \end{beamercolorbox}}%
  \vskip1pt
}

% full-bleed rust title page
\setbeamertemplate{title page}{%
  \begin{tikzpicture}[remember picture, overlay]
    \fill[ThemeRust] (current page.south west) rectangle (current page.north east);
  \end{tikzpicture}%
  \vfill
  \begin{flushright}
    {\headfont\bfseries\LARGE\color{ThemeCream}\inserttitle\par}
    \vspace{0.3em}
    {\color{ThemeCream!85}\small\insertsubtitle\par}
    \vspace{1.3em}
    {\color{ThemeCream}\small\insertauthor\par}
    \vspace{0.5em}
    {\color{ThemeCream!80}\scriptsize\insertinstitute\par}
    \vspace{0.8em}
    {\color{ThemeCream!75}\scriptsize\insertdate\par}
  \end{flushright}
  \vfill
}

% --- Deck palette remapped onto the theme -------------------------------
% Overrides the navy/violet set defined at the top of this file. Every TikZ
% figure and tcolorbox references these names, so they all follow along.
\definecolor{PrimaryColor}{HTML}{9B4632}     % rust: highlights, frames, titles
\definecolor{SecondaryColor}{HTML}{7A5C4B}   % warm brown: secondary boxes
\definecolor{AccentColor}{HTML}{7E8C57}      % sage: arrows, alerted text
\definecolor{ObjectColor}{HTML}{A9745B}      % clay: solid shapes
\definecolor{LightColor}{HTML}{D9A441}       % warm gold
\definecolor{TextColor}{HTML}{191919}
\definecolor{BackgroundColor}{HTML}{F8F6F2}
\definecolor{LightGray}{HTML}{E7E7DB}
\definecolor{DarkGray}{HTML}{6B6257}

% the four fern rules -- must match COLS in part 3/figs5.py
\definecolor{RuleA}{HTML}{C0392B}
\definecolor{RuleB}{HTML}{2E7D32}
\definecolor{RuleC}{HTML}{B8860B}
\definecolor{RuleD}{HTML}{1F6F8B}



\title{Fractals}
\subtitle{Infinite Details from Simple Rules}
\author{Speaker 1, Speaker 2, Speaker 3}   % <-- names
\institute{Department of Computer Science and Engineering\\
  Bangladesh University of Engineering and Technology (BUET)}
\date{CSE 200 Presentation}

% one-line takeaway box used on every slide
\newcommand{\takeaway}[1]{%
  \begin{tcolorbox}[colback=PrimaryColor!8, colframe=PrimaryColor, arc=2mm,
    boxrule=0.8pt, left=6pt, right=6pt, top=3pt, bottom=3pt]
    \small #1
  \end{tcolorbox}}
% small muted source line
\newcommand{\source}[1]{{\scriptsize\color{DarkGray}#1}}


\begin{document}

\begin{frame}
  \titlepage
\end{frame}

% ================================================================
% PART 1 — HISTORY (3 frames): the crisis, the pre-Mandelbrot
% "monsters", and Mandelbrot's reframing.
% ================================================================
\section{History}

% ---------------------------------------------------------------
% Frame 1 of 3: The crisis — an assumption calculus took for granted
% ---------------------------------------------------------------
\begin{frame}{A Crisis in 19th-Century Mathematics}
  \small
  For centuries, calculus rested on an assumption that felt obvious:

  \begin{conceptbox}{The Assumption}
    Any continuous curve is \highlight{smooth almost everywhere} —
    zoom in far enough on any point, and the curve should straighten
    out into a well-defined tangent line.
  \end{conceptbox}

  \pause
  \vspace{0.3cm}
  This wasn't just intuition — it was treated as close to self-evident.
  \highlight{Continuity} was assumed to imply \highlight{differentiability},
  almost everywhere, almost always.

  \pause
  \vspace{0.3cm}
  \begin{center}
    \begin{tikzpicture}
      \node[conceptbox, text width=0.75\textwidth] {
        \centering
        In 1872, \highlight{Karl Weierstrass} broke this assumption —
        on purpose.
      };
    \end{tikzpicture}
  \end{center}

  \pause
  \vspace{0.2cm}
  He constructed a function that is \highlight{continuous everywhere},
  yet \highlight{differentiable nowhere}: infinitely jagged, at every
  single scale, with no tangent line anywhere at all.
\end{frame}

% ---------------------------------------------------------------
% Frame 2 of 3: The "monster" era — mathematicians before Mandelbrot
% ---------------------------------------------------------------
\begin{frame}{Mathematical "Monsters"}
  \small
  Weierstrass's function was treated as a \highlight{pathological
  monster} — proof that pictures and intuition could not be trusted,
  and that analysis needed rigorous proof instead.

  \vspace{0.2cm}
  \pause
  This set off a chain reaction. Mathematicians began deliberately
  constructing stranger and stranger shapes, to probe exactly where
  classical geometry broke down:

  \begin{itemize}
    \item<3-> \highlight{Georg Cantor} (1883) — the Cantor set: an
      infinite set of points with zero total length.
    \item<4-> \highlight{Helge von Koch} (1904) — a geometric,
      drawable curve that is nowhere smooth: the \highlight{Koch curve}.
    \item<5-> \highlight{Wac\l{}aw Sierpi\'nski} — self-similar shapes
      like the \highlight{Sierpi\'nski triangle}.
    \item<6-> \highlight{Felix Hausdorff} (1918) — gave these shapes a
      number: a \highlight{fractional dimension}, between the familiar
      integer dimensions.
    \item<7-> \highlight{Gaston Julia \& Pierre Fatou} ($\sim$1918) —
      studied iterated functions in the complex plane, without any way
      to actually see what they looked like.
  \end{itemize}

  \pause[8]
  \vspace{0.2cm}
  \begin{center}
    \begin{tikzpicture}
      \node[conceptbox, text width=0.8\textwidth] {
        \centering
        For nearly a century, these were considered
        \highlight{beautiful but useless} — curiosities built to
        stress-test the foundations of mathematics, not to describe
        anything real.
      };
    \end{tikzpicture}
  \end{center}
\end{frame}

% ---------------------------------------------------------------
% Frame 3 of 3: Mandelbrot's question and reframing
% ---------------------------------------------------------------
\begin{frame}{Mandelbrot's Question}
  \small
  In the 1960s, working at IBM, \highlight{Benoit Mandelbrot} asked a
  deceptively simple question:

  \begin{center}
    \begin{tikzpicture}
      \node[conceptbox, text width=0.8\textwidth] {
        \centering
        \textit{``How long is the coast of Britain?''}
      };
    \end{tikzpicture}
  \end{center}

  \pause
  \vspace{0.2cm}
  The honest answer: \highlight{it depends on your ruler}. Measure
  with a 100~km ruler, get one length. Use a 1~km ruler, and the
  measured length grows — every bay and inlet you'd smoothed over
  now counts. The length never converges to a single number.

  \pause
  \vspace{0.2cm}
  Mandelbrot's insight: this isn't a flaw in the coastline — it's a
  sign that classical length and area are \highlight{the wrong tools}
  for something like this. What a coastline actually needs is exactly
  the "useless" mathematics of Cantor, Koch, and Hausdorff.

  \pause
  \vspace{0.2cm}
  \begin{itemize}
    \item He had a real problem no one had connected to these ideas.
    \item He had \highlight{computers} — something Cantor, Koch, and
      Julia never had — letting him \emph{see} these shapes for the
      first time.
    \item He united a century of scattered ``monsters'' under one
      name: \highlight{fractal}, from the Latin \textit{fractus}
      (``broken'').
  \end{itemize}
\end{frame}

% ================================================================
% PART 2 — DEFINITION + SIERPINSKI CONSTRUCTION WALKTHROUGH
% ================================================================
\section{Definition}

\begin{frame}{What are Fractals?}
  \begin{conceptbox}{Definition}
    A \highlight{fractal} is a geometric shape containing detailed
    structure at arbitrarily small scales, usually having a
    \highlight{fractal dimension} strictly exceeding the
    \highlight{topological dimension}.
  \end{conceptbox}

  \vspace{0.1cm}

  \textbf{Key Characteristics:}
  \begin{itemize}
    \item \highlight{Scale invariance}: Details at every level of magnification
    \item \highlight{Fractal dimension}: More than the \textit{conventional} dimension
    \item \highlight{Infinite complexity}: Generated by simple, recursive rules
      \begin{itemize}
        \item Like the infinite complexities of nature
      \end{itemize}
  \end{itemize}

  \vspace{0.1cm}

  \pause
  \begin{center}
    \begin{tikzpicture}
      \node[conceptbox, text width=0.7\textwidth] {
        \centering
        \textit{"Clouds are not spheres, mountains are not cones, coastlines are not circles..."} \\
        \textbf{--- Benoit Mandelbrot}
      };
    \end{tikzpicture}
  \end{center}
\end{frame}

\section{Some Fractals}

\begin{frame}{The Sierpinski Triangle}
  \begin{columns}
    \begin{column}{0.6\textwidth}
      \small
      \only<1-4>{
        Let's construct the Sierpinski Triangle:
        \begin{conceptbox}{Construction Process}
          \footnotesize
          \begin{enumerate}
            \item Start with an equilateral triangle
            \item<2-> Remove the central triangle (triangle connecting midpoints of the sides)
            \item<3-> Repeat for each triangle
            \item<4-> After infinite iterations, we get the Sierpinski Triangle
          \end{enumerate}
        \end{conceptbox}
      }
    \end{column}
    \begin{column}{0.4\textwidth}
      \begin{center}
        \begin{figure}
          \foreach \t/\i in {1/0,2/1,3/2,4/3} {
            \only<\t>{
              \includegraphics[width=\textwidth]{images/sierpinski/\i.png}
              \caption*{Iteration \i}
            }
          }
        \end{figure}
      \end{center}
    \end{column}
  \end{columns}
\end{frame}

% ================================================================
% PART 3 — SIERPINSKI AREA & PERIMETER CALCULATION
% ================================================================
\section{Area \& Perimeter}

\begin{frame}{The Sierpinski Triangle: Area}
  \begin{columns}
    \begin{column}{0.6\textwidth}
      \small
      Now let's find the area of triangle(s):
      \begin{mathbox}{Area Calculation}
        \footnotesize
        \begin{align*}
          \only<1->{
            A_0 &= \frac{1}{2}\cdot \sin(60^\circ) \cdot a^2 = \frac{\sqrt{3}}{4} a^2 \\
          }
          \only<2>{
            A_1 &= A_0 - \frac{1}{2}\cdot \sin(60^\circ) \cdot \left(\frac{a}{2}\right)^2 \\
            &= A_0 - \frac{\sqrt{3}}{16} a^2 = A_0 - \frac{1}{4} A_0 \\
            &= \frac{3}{4} A_0
          }
          \only<3->{
            A_1  &= \left(1-\frac{1}{4}\right) A_0 =  \frac{3}{4} A_0 \\
            A_2 &= \left(1-\frac{1}{4}\right) A_1 = \left(\frac{3}{4}\right)^2 A_0 \\
          }
          \only<4->{
            \hdots \\
            A_n &= \left(1-\frac{1}{4}\right)^n A_0 = \left(\frac{3}{4}\right)^n A_0 \\
            \lim_{n \to \infty} A_n &= 0
          }
        \end{align*}
      \end{mathbox}
      \only<2-3>{
        We cut out one fourth of the area at each step. So, naturally
        area is reduced by $\frac{3}{4}$ at each step.
      }
      \only<4->{
        \highlight{The area of the Sierpinski Triangle is zero even
        though it is a 2D shape.}
      }
    \end{column}
    \begin{column}{0.4\textwidth}
      \begin{center}
        \begin{figure}
          \foreach \t/\i in {1/0, 2/1, 3/2, 4/3} {
            \only<\t>{
              \includegraphics[width=\textwidth]{images/sierpinski/\i.png}
              \caption*{Iteration \i}
            }
          }
        \end{figure}
      \end{center}
    \end{column}
  \end{columns}
\end{frame}

\begin{frame}{The Sierpinski Triangle: Perimeter}
  \begin{columns}
    \begin{column}{0.6\textwidth}
      \small
      Now let's find the perimeter:
      \begin{mathbox}{Perimeter Calculation}
        \footnotesize
        \begin{align*}
          \only<1->{
            P_0 &= 3a \\
          }
          \only<2->{
            P_1 &= P_0 + 3\cdot\frac{a}{2} = \frac{9}{2} a = \frac{3}{2} P_0 \\
          }
          \only<3->{
            P_2 &= P_1 + 3\cdot3\cdot\frac{a}{4} = P_1 + \frac{1}{2} P_1 = \frac{3}{2} P_1 \\
          }
          \only<4->{
            \hdots \\
            P_n &= \frac{3}{2} P_{n-1} = \left(\frac{3}{2}\right)^n P_0 \\
            \lim_{n \to \infty} P_n &= \infty
          }
        \end{align*}
      \end{mathbox}
      \only<2-3>{
        We add half of the previous perimeter at each step.
        This is because we add the same number of triangles as the
        last step, but each triangle is half the size of the
        previous one.

        As a result, the perimeter grows by a factor of $\frac{3}{2}$
        at each step.
      }
      \only<4->{
        \highlight{The Sierpinski Triangle has infinite perimeter but
        zero area.}
      }
    \end{column}
    \begin{column}{0.4\textwidth}
      \begin{center}
        \begin{figure}
          \foreach \t/\i in {1/0, 2/1, 3/2, 4/3} {
            \only<\t>{
              \includegraphics[width=\textwidth]{images/sierpinski/\i.png}
              \caption*{Iteration \i}
            }
          }
        \end{figure}
      \end{center}
    \end{column}
  \end{columns}
\end{frame}

% ================================================================
% Handoff from Speaker 1
% ================================================================
\begin{frame}{Handoff to Speaker 2}
  \small
  We've now seen something strange, and precise:

  \begin{itemize}
    \item A \highlight{2D shape} (Sierpinski Triangle) with
      \highlight{zero area}.
    \item A shape built from a rule so simple a child could apply it
      by hand.
    \item No matter how far we zoom in, there's always more detail.
  \end{itemize}

  \pause
  \vspace{0.3cm}
  \begin{center}
    \begin{tikzpicture}
      \node[conceptbox, text width=0.8\textwidth] {
        \centering
        If it isn't a line, and it isn't a surface --- what
        \highlight{is} its true dimension?
      };
    \end{tikzpicture}
  \end{center}

  \vspace{0.2cm}
  \pause
  That number is exactly what Hausdorff was reaching for in 1918 ---
  and it's where Speaker~2 picks up the story.
\end{frame}

% ================================================================
% PART 2 — FRACTAL DIMENSION (Speaker 2)
% 9 frames: transition, integer dims, Sierpinski reveal,
%   scaling, formula, Sierpinski derivation, Koch derivation,
%   box-counting theory, handoff to Speaker 3
% ================================================================
\section{Fractal Dimension}

% ---------------------------------------------------------------
% S2 Frame 1: Transition — section title
% ---------------------------------------------------------------
\begin{frame}{}
  \vfill
  \begin{center}
    {\large\color{SecondaryColor}\itshape
      Let's talk about a property of fractals\ldots}

    \vspace{0.8cm}
    \begin{tikzpicture}
      \draw[PrimaryColor, line width=1.5pt] (-3.5,0) -- (3.5,0);
      \foreach \x in {-3,-2.5,...,3}
        \draw[AccentColor, line width=0.8pt] (\x,0.12) -- (\x,-0.12);
    \end{tikzpicture}

    \vspace{0.8cm}
    {\fontsize{26}{32}\selectfont\bfseries\color{PrimaryColor}
      Fractal Dimension}

    \vspace{0.5cm}
    {\small\color{DarkGray}
      What does it mean for a shape to have a
      \highlight{non-integer} dimension?}
  \end{center}
  \vfill
\end{frame}

% ---------------------------------------------------------------
% S2 Frame 2: Dimensions We Know (3 cumulative overlays)
% ---------------------------------------------------------------
\begin{frame}{Dimensions We Know}
  \small
  We are all familiar with integer dimensions:
  \vspace{0.2cm}
  \begin{center}
    \begin{tikzpicture}[
        dbox/.style={rounded corners=3pt, draw=PrimaryColor,
          fill=PrimaryColor!12, inner sep=5pt,
          font=\Large\bfseries, text=PrimaryColor}
      ]
      % ── Row 1: Line → 1D ──
      \onslide<1->{
        \draw[ObjectColor, line width=3pt] (-4,2) -- (-1,2);
        \node[font=\footnotesize, text=DarkGray, below]
          at (-2.5,1.7) {Line};
        \draw[-{Stealth[length=5pt]}, thick, AccentColor]
          (-0.3,2) -- (0.5,2);
        \node[dbox] at (1.6,2) {1D};
      }

      % ── Row 2: Square → 2D ──
      \onslide<2->{
        \draw[ObjectColor, line width=1.5pt, fill=ObjectColor!18]
          (-4,-0.8) rectangle (-1.6,0.8);
        \node[font=\footnotesize, text=DarkGray, below]
          at (-2.8,-1) {Square};
        \draw[-{Stealth[length=5pt]}, thick, AccentColor]
          (-0.3,0) -- (0.5,0);
        \node[dbox] at (1.6,0) {2D};
      }

      % ── Row 3: Cube → 3D ──
      \onslide<3->{
        \begin{scope}[shift={(-4,-3.5)}]
          % front face
          \draw[ObjectColor, line width=1.5pt, fill=ObjectColor!20]
            (0,0) rectangle (1.8,1.8);
          % top face
          \draw[ObjectColor, line width=1.5pt, fill=ObjectColor!28]
            (0,1.8) -- (0.55,2.3) -- (2.35,2.3)
            -- (1.8,1.8) -- cycle;
          % right face
          \draw[ObjectColor, line width=1.5pt, fill=ObjectColor!12]
            (1.8,0) -- (2.35,0.5) -- (2.35,2.3)
            -- (1.8,1.8) -- cycle;
        \end{scope}
        \node[font=\footnotesize, text=DarkGray, below]
          at (-2.8,-3.75) {Cube};
        \draw[-{Stealth[length=5pt]}, thick, AccentColor]
          (-0.3,-2.6) -- (0.5,-2.6);
        \node[dbox] at (1.6,-2.6) {3D};
      }
    \end{tikzpicture}
  \end{center}

  \onslide<3->{
    \begin{center}
      {\itshape\color{DarkGray}
        Dimensions are always whole numbers\ldots\ right?}
    \end{center}
  }
\end{frame}

% ---------------------------------------------------------------
% S2 Frame 3: What About This One? (2 overlays)
% ---------------------------------------------------------------
\begin{frame}{What About This One?}
  \begin{columns}
    \begin{column}{0.55\textwidth}
      \only<1>{
        \small
        \vspace{0.5cm}
        {\large What dimension does this shape have?}

        \vspace{0.4cm}
        {\itshape Any guesses?}

        \vspace{0.3cm}
        \begin{itemize}
          \item It lives in 2D space\ldots
          \item But it has \highlight{zero area}
          \item And \highlight{infinite perimeter}
        \end{itemize}

        \vspace{0.3cm}
        It's not quite 2D, but it's clearly more than 1D\ldots
      }
      \only<2>{
        \begin{conceptbox}{Sierpinski Dimension}
          The Sierpinski Triangle has a fractal dimension of:
          \begin{center}
            {\fontsize{28}{34}\selectfont
              $D \approx \highlight{1.585}$}

            \vspace{0.1cm}
            {\Large\color{RayColor}\textbf{!!}}
          \end{center}

          \vspace{0.2cm}
          A \highlight{non-integer} dimension --- between a line
          (1D) and a filled surface (2D).
        \end{conceptbox}
      }
    \end{column}
    \begin{column}{0.4\textwidth}
      \begin{center}
        \begin{figure}
          \includegraphics[width=0.85\textwidth]
            {images/sierpinski/5.png}
          \only<1>{\caption*{Sierpinski Triangle}}
          \only<2>{\caption*{$D=\log3/\log2\approx1.585$}}
        \end{figure}
      \end{center}
    \end{column}
  \end{columns}
\end{frame}

% ---------------------------------------------------------------
% S2 Frame 4: Self-Similarity & Scaling (5 overlays)
% ---------------------------------------------------------------
\begin{frame}{Self-Similarity \& Scaling}
  \centering

  % ── Overlay 1: Scaling a Line ──
  \only<1>{%
    \begin{tcolorbox}[colback=PrimaryColor!8, colframe=PrimaryColor, arc=2mm, boxrule=1pt,
                      width=0.96\textwidth, halign=center, top=3pt, bottom=3pt]
      \textbf{\color{PrimaryColor}Scaling a Line ($D = 1$)}:\quad
      Scale by $r = \frac{1}{2} \implies$ \highlight{2 pieces}.
      \quad Scale again by $\frac{1}{2} \implies 2 \times 2 = \mathbf{4\text{ pieces}}$.
    \end{tcolorbox}

    \vspace{0.45cm}
    \begin{center}
      \begin{tikzpicture}[scale=0.88]
        % Original Line
        \draw[ObjectColor, line width=3.5pt] (-6.0, 0) -- (-3.4, 0);
        \node[font=\footnotesize\bfseries, PrimaryColor, below=4pt]
          at (-4.7, 0) {Original Line};

        % Arrow 1
        \draw[-{Stealth[scale=1.2]}, line width=2pt, AccentColor]
          (-2.9, 0) -- (-1.9, 0);
        \node[font=\small\bfseries, above=2pt, AccentColor]
          at (-2.4, 0.1) {$\times\frac{1}{2}$};

        % 2 Pieces
        \draw[ObjectColor, line width=3.5pt] (-1.4, 0) -- (-0.2, 0);
        \draw[ObjectColor, line width=3.5pt] (0.1, 0) -- (1.3, 0);
        \node[font=\footnotesize\bfseries, PrimaryColor, below=4pt]
          at (-0.05, 0) {\textbf{2} pieces};

        % Arrow 2
        \draw[-{Stealth[scale=1.2]}, line width=2pt, AccentColor]
          (1.8, 0) -- (2.8, 0);
        \node[font=\small\bfseries, above=2pt, AccentColor]
          at (2.3, 0.1) {$\times\frac{1}{2}$};

        % 4 Pieces
        \foreach \k in {0, 1, 2, 3} {
          \draw[ObjectColor, line width=3.5pt]
            (3.3 + \k*0.75, 0) -- ++(0.55, 0);
        }
        \node[font=\footnotesize\bfseries, PrimaryColor, below=4pt]
          at (4.7, 0) {\textbf{4} pieces};
      \end{tikzpicture}
    \end{center}
  }%

  % ── Overlay 2: Scaling a Square — Step 1 (4 pieces) ──
  \only<2>{%
    \begin{tcolorbox}[colback=PrimaryColor!8, colframe=PrimaryColor, arc=2mm, boxrule=1pt,
                      width=0.96\textwidth, halign=center, top=3pt, bottom=3pt]
      \textbf{\color{PrimaryColor}Scaling a Square ($D = 2$) --- Step 1}:\quad
      Scale by $r = \frac{1}{2} \implies$ \highlight{4 self-similar pieces} ($2 \times 2$ grid).
    \end{tcolorbox}

    \vspace{0.45cm}
    \begin{center}
      \begin{tikzpicture}[scale=0.85]
        % Original Square
        \draw[ObjectColor, line width=1.5pt, fill=ObjectColor!18]
          (-5.5, -0.9) rectangle ++(1.8, 1.8);
        \node[font=\footnotesize\bfseries, PrimaryColor, below=4pt]
          at (-4.6, -0.9) {Original Square};

        % Arrow 1
        \draw[-{Stealth[scale=1.1]}, line width=1.8pt, AccentColor]
          (-3.2, 0) -- (-2.2, 0);
        \node[font=\scriptsize\bfseries, above=1pt, AccentColor]
          at (-2.7, 0.1) {$\times\frac{1}{2}$};

        % 4 Pieces (2x2 grid)
        \foreach \i in {0, 1} {
          \foreach \j in {0, 1} {
            \draw[ObjectColor, line width=1pt, fill=ObjectColor!18]
              (-1.8 + \i*1.05, -0.85 + \j*1.05) rectangle ++(0.8, 0.8);
          }
        }
        \node[font=\footnotesize\bfseries, PrimaryColor, below=4pt]
          at (-0.95, -0.85) {\textbf{4} pieces};

        % Result Badge
        \node[font=\normalsize\bfseries, PrimaryColor, right]
          at (0.7, 0) {$N = 2^2 = \mathbf{4}$};
      \end{tikzpicture}
    \end{center}
  }%

  % ── Overlay 3: Scaling a Square — Step 2 (16 pieces) ──
  \only<3>{%
    \begin{tcolorbox}[colback=PrimaryColor!8, colframe=PrimaryColor, arc=2mm, boxrule=1pt,
                      width=0.96\textwidth, halign=center, top=3pt, bottom=3pt]
      \textbf{\color{PrimaryColor}Scaling a Square ($D = 2$) --- Step 2}:\quad
      Scale each piece by $\frac{1}{2}$ again $\implies 4 \times 4 = \mathbf{16\text{ self-similar pieces}}$ ($4 \times 4$ grid)!
    \end{tcolorbox}

    \vspace{0.45cm}
    \begin{center}
      \begin{tikzpicture}[scale=0.85]
        % Stage 1: Original Square
        \draw[ObjectColor, line width=1.5pt, fill=ObjectColor!18]
          (-5.5, -0.9) rectangle ++(1.8, 1.8);
        \node[font=\footnotesize\bfseries, PrimaryColor, below=4pt]
          at (-4.6, -0.9) {Original};

        % Arrow 1
        \draw[-{Stealth[scale=1.1]}, line width=1.8pt, AccentColor]
          (-3.2, 0) -- (-2.2, 0);
        \node[font=\scriptsize\bfseries, above=1pt, AccentColor]
          at (-2.7, 0.1) {$\times\frac{1}{2}$};

        % Stage 2: 4 Pieces (2x2 grid)
        \foreach \i in {0, 1} {
          \foreach \j in {0, 1} {
            \draw[ObjectColor, line width=1pt, fill=ObjectColor!18]
              (-1.8 + \i*1.05, -0.85 + \j*1.05) rectangle ++(0.8, 0.8);
          }
        }
        \node[font=\footnotesize\bfseries, PrimaryColor, below=4pt]
          at (-0.95, -0.85) {\textbf{4} pieces};

        % Arrow 2
        \draw[-{Stealth[scale=1.1]}, line width=1.8pt, AccentColor]
          (0.6, 0) -- (1.6, 0);
        \node[font=\scriptsize\bfseries, above=1pt, AccentColor]
          at (1.1, 0.1) {$\times\frac{1}{2}$};

        % Stage 3: 16 Pieces (4 clusters of 2x2 mini-squares)
        \foreach \ci in {0, 1} {
          \foreach \cj in {0, 1} {
            \foreach \i in {0, 1} {
              \foreach \j in {0, 1} {
                \draw[ObjectColor, line width=0.8pt, fill=ObjectColor!18]
                  (2.1 + \ci*1.02 + \i*0.44, -0.92 + \cj*1.02 + \j*0.44) rectangle ++(0.36, 0.36);
              }
            }
          }
        }
        \node[font=\footnotesize\bfseries, PrimaryColor, below=4pt]
          at (2.95, -0.92) {\textbf{16} pieces};

        % Result Badge
        \node[font=\normalsize\bfseries, PrimaryColor, right]
          at (4.5, 0) {$N = 4^2 = \mathbf{16}$};
      \end{tikzpicture}
    \end{center}
  }%

  % ── Overlay 4: Scaling a Cube ──
  \only<4>{%
    \begin{tcolorbox}[colback=PrimaryColor!8, colframe=PrimaryColor, arc=2mm, boxrule=1pt,
                      width=0.96\textwidth, halign=center, top=3pt, bottom=3pt]
      \textbf{\color{PrimaryColor}Scaling a Cube ($D = 3$)}:\quad
      Scale by $r = \frac{1}{2} \implies$ \highlight{8 self-similar pieces} ($2 \times 2 \times 2$ grid in 3D).
    \end{tcolorbox}

    \vspace{0.35cm}
    \begin{center}
      \begin{tikzpicture}[scale=0.80]
        % Left: Original Cube
        \filldraw[draw=ObjectColor, line width=1.5pt, fill=ObjectColor!30]
          (-4.2, 0.9) -- ++(0.65, 0.65) -- ++(1.8, 0) -- ++(-0.65, -0.65) -- cycle;
        \filldraw[draw=ObjectColor, line width=1.5pt, fill=ObjectColor!20]
          (-2.4, -0.9) -- ++(0.65, 0.65) -- ++(0, 1.8) -- ++(-0.65, -0.65) -- cycle;
        \filldraw[draw=ObjectColor, line width=1.5pt, fill=ObjectColor!14]
          (-4.2, -0.9) rectangle ++(1.8, 1.8);
        \node[font=\footnotesize\bfseries, PrimaryColor, below=4pt]
          at (-2.95, -1.1) {Original Cube};

        % Center: Transition Arrow
        \draw[-{Stealth[scale=1.2]}, line width=2pt, AccentColor]
          (-1.4, 0.1) -- (-0.1, 0.1);
        \node[font=\small\bfseries, above=2pt, AccentColor]
          at (-0.75, 0.25) {scale $\times\frac{1}{2}$};

        % Right: 8 Separated Small Cubes
        % Back layer (Z = 1, offset by (0.85, 0.85))
        \foreach \i in {0, 1} {
          \foreach \j in {0, 1} {
            \begin{scope}[shift={(0.6 + 0.85 + \i*1.30, -1.0 + 0.85 + \j*1.30)}]
              \filldraw[draw=ObjectColor, line width=1pt, fill=ObjectColor!30]
                (0, 0.75) -- ++(0.35, 0.35) -- ++(0.75, 0) -- ++(-0.35, -0.35) -- cycle;
              \filldraw[draw=ObjectColor, line width=1pt, fill=ObjectColor!20]
                (0.75, 0) -- ++(0.35, 0.35) -- ++(0, 0.75) -- ++(-0.35, -0.35) -- cycle;
              \filldraw[draw=ObjectColor, line width=1pt, fill=ObjectColor!14]
                (0, 0) rectangle ++(0.75, 0.75);
            \end{scope}
          }
        }

        % Front layer (Z = 0)
        \foreach \i in {0, 1} {
          \foreach \j in {0, 1} {
            \begin{scope}[shift={(0.6 + \i*1.30, -1.0 + \j*1.30)}]
              \filldraw[draw=ObjectColor, line width=1pt, fill=ObjectColor!30]
                (0, 0.75) -- ++(0.35, 0.35) -- ++(0.75, 0) -- ++(-0.35, -0.35) -- cycle;
              \filldraw[draw=ObjectColor, line width=1pt, fill=ObjectColor!20]
                (0.75, 0) -- ++(0.35, 0.35) -- ++(0, 0.75) -- ++(-0.35, -0.35) -- cycle;
              \filldraw[draw=ObjectColor, line width=1pt, fill=ObjectColor!14]
                (0, 0) rectangle ++(0.75, 0.75);
            \end{scope}
          }
        }

        \node[font=\footnotesize\bfseries, PrimaryColor, below=4pt]
          at (2.3, -1.1) {\textbf{8} separated cubes};

        % Result Badge
        \node[font=\normalsize\bfseries, PrimaryColor, right]
          at (4.3, 0.1) {$N = 2^3 = \mathbf{8}$};
      \end{tikzpicture}
    \end{center}
  }%

  % ── Overlay 5: Do You See a Pattern? ──
  \only<5>{%
    \begin{tcolorbox}[colback=PrimaryColor!8, colframe=PrimaryColor, arc=2mm, boxrule=1pt,
                      width=0.96\textwidth, halign=center, top=3pt, bottom=3pt]
      \textbf{\color{PrimaryColor}Do You See a Pattern?}\\
      Each scaling step multiplies pieces by $(1/r)^D$!
    \end{tcolorbox}

    \vspace{0.25cm}
    \begin{columns}
      \begin{column}{0.55\textwidth}
        \centering
        \small
        \renewcommand{\arraystretch}{1.3}
        \begin{tabular}{lcccc}
          \toprule
          \textbf{Shape} & \textbf{Scale} & \textbf{1st Step} & \textbf{2nd Step} & \textbf{$D$} \\
          \midrule
          Line   & $1/2$ & $2 = 2^1$ & $4 = 2^2$  & \textbf{1} \\
          Square & $1/2$ & $4 = 2^2$ & $16 = 2^4$ & \textbf{2} \\
          Cube   & $1/2$ & $8 = 2^3$ & $(64 = 2^6)$ & \textbf{3} \\
          \bottomrule
        \end{tabular}

        \vspace{0.3cm}
        \begin{conceptbox}{The Pattern}
          \centering
          \highlight{Number of pieces} $= \left(\dfrac{1}{r}\right)^{\!D}$
        \end{conceptbox}
      \end{column}
      \begin{column}{0.4\textwidth}
        \begin{center}
          \includegraphics[width=0.45\textwidth]{images/sierpinski/3.png}

          \vspace{0.2cm}
          {\Large \textbf{\color{PrimaryColor}$D = \;?$}}\\
          \vspace{0.1cm}
          \footnotesize What if the shape is a \highlight{fractal}?
        \end{center}
      \end{column}
    \end{columns}
  }%
\end{frame}

% ---------------------------------------------------------------
% S2 Frame 5: The Self-Similarity Rule (2 overlays)
% ---------------------------------------------------------------
\begin{frame}{The Self-Similarity Rule}
  \centering
  \small
  After scaling by factor $r$, the number of self-similar pieces is:
  \quad
  \highlight{\Large $N = \left(\dfrac{1}{r}\right)^{\!D}$}

  \vspace{0.25cm}
  \only<1>{%
    {\normalsize \textbf{\color{PrimaryColor}Square ($D = 2$)} \quad with scale factor $r = \frac{1}{2}$:\quad
    $N = \left(\frac{1}{1/2}\right)^{\!2} = 2^2 = \mathbf{4\text{ pieces}}$}%
  }%
  \only<2>{%
    {\normalsize \textbf{\color{PrimaryColor}Cube ($D = 3$)} \quad with scale factor $r = \frac{1}{2}$:\quad
    $N = \left(\frac{1}{1/2}\right)^{\!3} = 2^3 = \mathbf{8\text{ pieces}}$}%
  }%

  \vspace{0.3cm}
  \only<1>{%
    \begin{center}
      \begin{tikzpicture}[scale=0.85]
        % ── Left: Original Square ──
        \draw[ObjectColor, line width=1.5pt, fill=ObjectColor!18]
          (-4.2, -1.0) rectangle ++(2.0, 2.0);
        \node[font=\footnotesize\bfseries, PrimaryColor, below]
          at (-3.2, -1.25) {Original Square};

        % ── Center: Transition Arrow ──
        \draw[-{Stealth[scale=1.2]}, line width=2pt, AccentColor]
          (-1.4, 0) -- (-0.1, 0);
        \node[font=\small\bfseries, above=2pt, AccentColor]
          at (-0.75, 0.15) {scale $\times\frac{1}{2}$};

        % ── Right: 4 Separated Squares ──
        \foreach \i in {0, 1} {
          \foreach \j in {0, 1} {
            \draw[ObjectColor, line width=1.2pt, fill=ObjectColor!18]
              (0.8 + \i*1.25, -1.075 + \j*1.25) rectangle ++(0.9, 0.9);
          }
        }
        \node[font=\footnotesize\bfseries, PrimaryColor, below]
          at (1.875, -1.25) {\textbf{4} separated pieces};

        % ── Result Badge ──
        \node[font=\normalsize\bfseries, PrimaryColor, right]
          at (3.6, 0) {$N = 2^2 = \mathbf{4}$};
      \end{tikzpicture}
    \end{center}
  }%
  \only<2>{%
    \begin{center}
      \begin{tikzpicture}[scale=0.80]
        % ── Left: Original Cube ──
        \filldraw[draw=ObjectColor, line width=1.5pt, fill=ObjectColor!30]
          (-4.2, 0.9) -- ++(0.65, 0.65) -- ++(1.8, 0) -- ++(-0.65, -0.65) -- cycle;
        \filldraw[draw=ObjectColor, line width=1.5pt, fill=ObjectColor!20]
          (-2.4, -0.9) -- ++(0.65, 0.65) -- ++(0, 1.8) -- ++(-0.65, -0.65) -- cycle;
        \filldraw[draw=ObjectColor, line width=1.5pt, fill=ObjectColor!14]
          (-4.2, -0.9) rectangle ++(1.8, 1.8);
        \node[font=\footnotesize\bfseries, PrimaryColor, below]
          at (-2.95, -1.25) {Original Cube};

        % ── Center: Transition Arrow ──
        \draw[-{Stealth[scale=1.2]}, line width=2pt, AccentColor]
          (-1.4, 0.1) -- (-0.1, 0.1);
        \node[font=\small\bfseries, above=2pt, AccentColor]
          at (-0.75, 0.25) {scale $\times\frac{1}{2}$};

        % ── Right: 8 Separated Small Cubes ──
        % Back layer (Z = 1, offset by (0.85, 0.85))
        \foreach \i in {0, 1} {
          \foreach \j in {0, 1} {
            \begin{scope}[shift={(0.6 + 0.85 + \i*1.30, -1.0 + 0.85 + \j*1.30)}]
              \filldraw[draw=ObjectColor, line width=1pt, fill=ObjectColor!30]
                (0, 0.75) -- ++(0.35, 0.35) -- ++(0.75, 0) -- ++(-0.35, -0.35) -- cycle;
              \filldraw[draw=ObjectColor, line width=1pt, fill=ObjectColor!20]
                (0.75, 0) -- ++(0.35, 0.35) -- ++(0, 0.75) -- ++(-0.35, -0.35) -- cycle;
              \filldraw[draw=ObjectColor, line width=1pt, fill=ObjectColor!14]
                (0, 0) rectangle ++(0.75, 0.75);
            \end{scope}
          }
        }

        % Front layer (Z = 0)
        \foreach \i in {0, 1} {
          \foreach \j in {0, 1} {
            \begin{scope}[shift={(0.6 + \i*1.30, -1.0 + \j*1.30)}]
              \filldraw[draw=ObjectColor, line width=1pt, fill=ObjectColor!30]
                (0, 0.75) -- ++(0.35, 0.35) -- ++(0.75, 0) -- ++(-0.35, -0.35) -- cycle;
              \filldraw[draw=ObjectColor, line width=1pt, fill=ObjectColor!20]
                (0.75, 0) -- ++(0.35, 0.35) -- ++(0, 0.75) -- ++(-0.35, -0.35) -- cycle;
              \filldraw[draw=ObjectColor, line width=1pt, fill=ObjectColor!14]
                (0, 0) rectangle ++(0.75, 0.75);
            \end{scope}
          }
        }

        \node[font=\footnotesize\bfseries, PrimaryColor, below]
          at (2.3, -1.25) {\textbf{8} separated small cubes};

        % ── Result Badge ──
        \node[font=\normalsize\bfseries, PrimaryColor, right]
          at (4.3, 0.1) {$N = 2^3 = \mathbf{8}$};
      \end{tikzpicture}
    \end{center}

    \vspace{0.15cm}
    \centering\small
    \textit{Can we use this same rule to find the dimension of a \highlight{fractal}?}
  }%
\end{frame}

% ---------------------------------------------------------------
% S2 Frame 6: Dimension of the Sierpinski Triangle (4 overlays)
% ---------------------------------------------------------------
\begin{frame}{Dimension of the Sierpinski Triangle}
  \begin{columns}
    \begin{column}{0.55\textwidth}
      \small
      \only<1>{
        Let's apply the self-similarity rule to the
        \highlight{Sierpinski Triangle}.

        \vspace{0.3cm}
        \begin{conceptbox}{Step 1}
          Scale the Sierpinski Triangle by $\frac{1}{2}$.
        \end{conceptbox}
      }
      \only<2>{
        \begin{conceptbox}{Step 2 --- Self-Similar Pieces}
          We get exactly \highlight{3} self-similar copies
          of the original triangle.

          \vspace{0.2cm}
          Each smaller triangle is a perfect replica,
          scaled by $r = \frac{1}{2}$.
        \end{conceptbox}
      }
      \only<3>{
        \begin{mathbox}{Step 3 --- The Equation}
          Using $N = \left(\frac{1}{r}\right)^D$:
          \begin{align*}
            3 &= \left(\frac{1}{1/2}\right)^{\!D} \\[6pt]
            3 &= 2^D
          \end{align*}
        \end{mathbox}
      }
      \only<4>{
        \begin{mathbox}{Step 4 --- Solving for $D$}
          \begin{align*}
            2^D &= 3 \\[4pt]
            D\,\log 2 &= \log 3 \\[4pt]
            D &= \frac{\log 3}{\log 2}
          \end{align*}
          \begin{center}
            {\fontsize{22}{28}\selectfont
              $D \approx \highlight{1.585}$}
          \end{center}
        \end{mathbox}

        \vspace{0.1cm}
        \highlight{Not an integer!} --- strictly between
        1D (line) and 2D (surface).
      }
    \end{column}
    \begin{column}{0.4\textwidth}
      \begin{center}
        \begin{tikzpicture}[scale=0.85]
          \only<1>{
            \node at (0,0) {
              \includegraphics[width=0.85\textwidth]
                {images/sierpinski/5.png}
            };
            \draw[-{Stealth}, very thick, AccentColor]
              (0,-1.8) -- (0,-2.5);
            \node[font=\small, AccentColor, right]
              at (0.1,-2.15) {scale by $\frac{1}{2}$};
          }
          \only<2-3>{
            % Three self-similar copies
            \node at (0,1.2) {
              \includegraphics[width=0.4\textwidth]
                {images/sierpinski/4.png}
            };
            \node at (-1,-0.6) {
              \includegraphics[width=0.4\textwidth]
                {images/sierpinski/4.png}
            };
            \node at (1,-0.6) {
              \includegraphics[width=0.4\textwidth]
                {images/sierpinski/4.png}
            };
            \node[font=\Large\bfseries, PrimaryColor]
              at (0,-2) {3 copies};
          }
          \only<4>{
            \node at (0,0.5) {
              \includegraphics[width=0.7\textwidth]
                {images/sierpinski/5.png}
            };
            \node[rounded corners=4pt, draw=PrimaryColor,
              fill=PrimaryColor!12, inner sep=6pt,
              font=\Large\bfseries, text=PrimaryColor]
              at (0,-2) {$D\approx1.585$};
          }
        \end{tikzpicture}
      \end{center}
    \end{column}
  \end{columns}
\end{frame}

% ---------------------------------------------------------------
% S2 Frame 7: Similarly — Koch Curve (2 overlays)
% ---------------------------------------------------------------
\begin{frame}{Similarly: The Koch Curve}
  \begin{columns}
    \begin{column}{0.55\textwidth}
      \small
      \only<1>{
        The \highlight{Koch Curve} can be divided into
        \highlight{4} self-similar pieces, each scaled by
        $r = \frac{1}{3}$.

        \vspace{0.3cm}
        \begin{conceptbox}{Setup}
          \begin{itemize}
            \item Scale factor: $r = \frac{1}{3}$
            \item Self-similar pieces: $N = 4$
          \end{itemize}
        \end{conceptbox}
      }
      \only<2>{
        \begin{mathbox}{Koch Dimension}
          \begin{align*}
            N &= \left(\frac{1}{r}\right)^{\!D} \\[4pt]
            4 &= 3^D \\[4pt]
            D &= \frac{\log 4}{\log 3}
          \end{align*}
          \begin{center}
            {\fontsize{22}{28}\selectfont
              $D \approx \highlight{1.262}$}
          \end{center}
        \end{mathbox}

        \vspace{0.2cm}
        Also \highlight{non-integer} --- between a line (1D)
        and a surface (2D), but closer to a line than
        the Sierpinski Triangle.
      }
    \end{column}
    \begin{column}{0.4\textwidth}
      \begin{center}
        \begin{tikzpicture}
          \only<1>{
            \node at (0,1.5) {
              \includegraphics[width=0.85\textwidth]
                {images/koch/5.png}
            };
            \draw[-{Stealth}, very thick, AccentColor]
              (0,0.3) -- (0,-0.4);
            \node[font=\small, AccentColor, right]
              at (0.1,-0.05) {$\times\frac{1}{3}$};
            % 4 smaller copies arranged in Koch shape
            \node at (-1.5,-1.5) {
              \includegraphics[width=0.22\textwidth]
                {images/koch/4.png}
            };
            \node[rotate=60] at (-0.45,-1.2) {
              \includegraphics[width=0.22\textwidth]
                {images/koch/4.png}
            };
            \node[rotate=-60] at (0.45,-1.2) {
              \includegraphics[width=0.22\textwidth]
                {images/koch/4.png}
            };
            \node at (1.5,-1.5) {
              \includegraphics[width=0.22\textwidth]
                {images/koch/4.png}
            };
            \node[font=\Large\bfseries, PrimaryColor]
              at (0,-2.8) {4 copies};
          }
          \only<2>{
            \node at (0,0.5) {
              \includegraphics[width=0.85\textwidth]
                {images/koch/5.png}
            };
            \node[rounded corners=4pt, draw=PrimaryColor,
              fill=PrimaryColor!12, inner sep=6pt,
              font=\Large\bfseries, text=PrimaryColor]
              at (0,-1.5) {$D\approx1.262$};
          }
        \end{tikzpicture}
      \end{center}
    \end{column}
  \end{columns}
\end{frame}

% ---------------------------------------------------------------
% S2 Frame 8: The General Theory — Box-Counting Dimension
% ---------------------------------------------------------------
\begin{frame}{The General Theory: Box-Counting Dimension}
  \small
  The \highlight{Self-Similarity Dimension} works beautifully for
  self-similar fractals. But what about more complex shapes?

  \vspace{0.3cm}
  \begin{mathbox}{Box-Counting Dimension}
    The fractal dimension $D$ is defined as:
    $$D = \lim_{\varepsilon \to 0}
      \frac{\log N(\varepsilon)}{\log(1/\varepsilon)}$$
    where $N(\varepsilon)$ is the number of boxes of side length
    $\varepsilon$ needed to cover the fractal.
  \end{mathbox}

  \vspace{0.3cm}
  \pause
  \begin{conceptbox}{Intuition}
    \begin{itemize}
      \item Cover the fractal with a grid of boxes of size
        $\varepsilon$
      \item Count how many boxes $N(\varepsilon)$ the fractal
        touches
      \item Shrink $\varepsilon \to 0$ and observe how
        $N(\varepsilon)$ grows
      \item The \highlight{rate of growth} tells us the
        dimension
    \end{itemize}
  \end{conceptbox}

  \vspace{0.2cm}
  \begin{center}
    {\footnotesize\color{DarkGray} For a deeper dive, see
    \href{https://www.youtube.com/watch?v=gB9n2gHsHN4}%
    {\textcolor{AccentColor}{3Blue1Brown:
    Fractals are typically not self-similar}}}
  \end{center}
\end{frame}

% ---------------------------------------------------------------
% S2 Frame 9: Handoff to Speaker 3
% ---------------------------------------------------------------
\begin{frame}{Handoff to Speaker 3}
  \small
  We've now uncovered a remarkable property of fractals:

  \begin{itemize}
    \item They have \highlight{non-integer dimensions} ---
      genuinely ``between'' the familiar 1D, 2D, 3D.
    \item The Sierpinski Triangle lives in dimension
      $D\approx 1.585$.
    \item The Koch Curve lives in dimension
      $D\approx 1.262$.
    \item The \highlight{box-counting dimension} generalises
      this to \emph{any} shape.
  \end{itemize}

  \pause
  \vspace{0.3cm}
  \begin{center}
    \begin{tikzpicture}
      \node[conceptbox, text width=0.8\textwidth] {
        \centering
        We now know \emph{what} fractals are and how to
        measure their complexity. \\[4pt]
        But \highlight{where do we see them}? And
        \highlight{how are they used} in computer science?
      };
    \end{tikzpicture}
  \end{center}

  \vspace{0.2cm}
  \pause
  That's exactly where Speaker~3 takes us ---
  fractals in nature, art, and computation.
\end{frame}

% ================================================================
% PART 3 — FRACTALS IN THE REAL WORLD (Speaker 3)
% 9 frames: section title, nature, coastline (2), body, database,
%   game worlds, limits, conclusion  (+ references)
% ================================================================
\section{Fractals in the Real World}

% ---------------------------------------------------------------
% S3 Frame 1: Transition — section title (same style as Speaker 2)
% ---------------------------------------------------------------
\begin{frame}{}
  \vfill
  \begin{center}
    {\large\color{SecondaryColor}\itshape
      We know what a fractal is, and how to measure it\ldots}

    \vspace{0.8cm}
    \begin{tikzpicture}
      \draw[PrimaryColor, line width=1.5pt] (-3.5,0) -- (3.5,0);
      \foreach \x in {-3,-2.5,...,3}
        \draw[AccentColor, line width=0.8pt] (\x,0.12) -- (\x,-0.12);
    \end{tikzpicture}

    \vspace{0.8cm}
    {\fontsize{26}{32}\selectfont\bfseries\color{PrimaryColor}
      Fractals in the Real World}

    \vspace{0.5cm}
    {\small\color{DarkGray}
      Where does all of this actually \highlight{show up}?}
  \end{center}
  \vfill
\end{frame}

% ---------------------------------------------------------------
% S3 Frame 2: Nature — real photos next to a rule-grown tree
% ---------------------------------------------------------------
\begin{frame}{Nature draws with one rule, repeated}
  \begin{columns}[T]
    \begin{column}{0.62\textwidth}
      \centering
      \includegraphics[width=0.42\textwidth]{images/part3/nature_romanesco.png}\hfill
      \includegraphics[width=0.42\textwidth]{images/part3/nature_snowflake.png}\\[3pt]
      \includegraphics[width=0.42\textwidth]{images/part3/nature_aloe.png}\hfill
      \includegraphics[width=0.42\textwidth]{images/part3/nature_pinecone.png}
    \end{column}
    \begin{column}{0.36\textwidth}
      \centering
      \includegraphics[width=0.75\textwidth]{images/tree/7.png}\\[2pt]
      {\footnotesize\color{DarkGray} grown by a computer from\\
        one rule: $F \rightarrow T[-F][+F]$}
    \end{column}
  \end{columns}
  \vspace{0.2cm}
  \takeaway{Broccoli, ice, leaves, cones: the \highlight{same shape at every zoom} ---
    exactly what the Sierpinski rule did on paper.}
\end{frame}

% ---------------------------------------------------------------
% S3 Frame 3: Coastline — Speaker 1's question, answered with data
% ---------------------------------------------------------------
\begin{frame}{A finer ruler makes Britain's coast longer}
  \centering
  \includegraphics[width=\textwidth]{images/part3/britain_rulers.png}

  \vspace{0.1cm}
  \takeaway{This is Speaker~1's question, on \highlight{real map data}:
    the finer the ruler, the more bays and inlets count.}

  \source{Coastline data: GSHHG (Wessel \& Smith 1996), plotted at three resolutions.}
\end{frame}

% ---------------------------------------------------------------
% S3 Frame 4: Richardson plot — Speaker 2's dimension, from data
% ---------------------------------------------------------------
\begin{frame}{Plot it on log--log axes: the slope is the dimension}
  \begin{columns}[c]
    \begin{column}{0.68\textwidth}
      \centering
      \includegraphics[width=\textwidth]{images/part3/richardson.png}
    \end{column}
    \begin{column}{0.31\textwidth}
      \begin{center}
        {\fontsize{26}{30}\selectfont\bfseries\color{PrimaryColor} $D \approx 1.3$}\\[6pt]
        {\small\color{DarkGray} Mandelbrot (1967)\\ found $1.25$}
      \end{center}
      \vspace{0.3cm}
      \takeaway{Speaker~2's box-counting idea, done with a ruler instead of boxes.}
    \end{column}
  \end{columns}
  \vspace{0.2cm}
  \source{Ruler-walk on the full-resolution GSHHG outline of Great Britain; rulers 4--256 km.}
\end{frame}

% ---------------------------------------------------------------
% S3 Frame 5: The body — branching packs area into volume
% ---------------------------------------------------------------
\begin{frame}{Your body branches to pack a huge surface into a small space}
  \begin{columns}[c]
    \begin{column}{0.58\textwidth}
      \centering
      \includegraphics[width=0.62\textwidth]{images/part3/retina.png}\hfill
      \raisebox{0.14\textwidth}{\includegraphics[width=0.34\textwidth]{images/part3/retina_zoom.png}}\\[2pt]
      {\footnotesize\color{DarkGray} blood vessels in a human eye, and a zoom-in}
    \end{column}
    \begin{column}{0.42\textwidth}
      \small
      \begin{itemize}
        \item Lungs: about \highlight{23 splits}\\
          $\Rightarrow$ $\sim$\highlight{70 m$^2$} of surface\\
          inside $\sim$5 litres
        \item Vessels: every cell within reach, using little space
        \item Doctors track the \highlight{dimension} of eye vessels to spot disease early
      \end{itemize}
    \end{column}
  \end{columns}
  \vspace{0.05cm}
  \takeaway{Branching again and again is nature's cheapest way to get \highlight{area out of volume}.}

  \source{Photo: M.~Häggström, ``Fundus photograph of normal left eye'', 2014, CC0. Lung figures: Weibel (1963).}
\end{frame}

% ---------------------------------------------------------------
% S3 Frame 6: Databases — space-filling curve = spatial index
% ---------------------------------------------------------------
\begin{frame}{A curve that fills a square turns a map into a sorted list}
  \begin{columns}[c]
    \begin{column}{0.5\textwidth}
      \centering
      \includegraphics[width=0.95\textwidth]{images/part3/hilbert_index.png}
    \end{column}
    \begin{column}{0.48\textwidth}
      \small
      \begin{itemize}
        \item A database can only sort by \highlight{one number}
        \item Walk the \highlight{Hilbert curve} through the map: each cell gets its position on the curve
        \item Nearby places $\rightarrow$ nearby numbers $\rightarrow$ one quick range search
      \end{itemize}
      \vspace{0.2cm}
      \takeaway{This trick powers map search in systems like Google~S2 and Uber~H3.}
    \end{column}
  \end{columns}
\end{frame}

% ---------------------------------------------------------------
% S3 Frame 7: Game worlds — grown from a rule, not stored
% ---------------------------------------------------------------
\begin{frame}{Game worlds are grown from a rule, not stored}
  \centering
  \includegraphics[width=0.98\textwidth]{images/part3/terrain_octaves.png}

  \vspace{0.15cm}
  \takeaway{Add \highlight{smaller bumps at every scale} --- the Koch idea --- and a
    handful of numbers becomes an endless mountain range. Store the rule, not the map.}

  \source{Terrain generated here by summing 7 layers of noise, each half the size of the last (fractal Brownian motion).}
\end{frame}

% ---------------------------------------------------------------
% S3 Frame 8: Limits — real fractals stop
% ---------------------------------------------------------------
\begin{frame}{Real fractals stop at some scale}
  \begin{columns}[c]
    \begin{column}{0.66\textwidth}
      \centering
      \includegraphics[width=\textwidth]{images/part3/richardson_limit.png}
    \end{column}
    \begin{column}{0.33\textwidth}
      \small
      \begin{itemize}
        \item Coast: bays look like bays, but \highlight{sand grains don't}
        \item Lungs: stop after $\sim$23 splits
        \item Only the Sierpinski triangle goes on \highlight{forever}
      \end{itemize}
      \vspace{0.2cm}
      \takeaway{A fractal is a \highlight{model} of nature --- and every model has a range where it works.}
    \end{column}
  \end{columns}
\end{frame}

% ---------------------------------------------------------------
% S3 Frame 9: Conclusion — the whole talk in one row
% ---------------------------------------------------------------
\begin{frame}{One rule, repeated, explains what smooth geometry couldn't}
  \begin{columns}[T]
    \begin{column}{0.33\textwidth}
      \centering
      {\small\bfseries\color{SecondaryColor} Part 1 --- the surprise}\\[4pt]
      \includegraphics[width=0.7\textwidth]{images/sierpinski/5.png}\\[4pt]
      {\small zero area,\\ infinite edge}
    \end{column}
    \begin{column}{0.33\textwidth}
      \centering
      {\small\bfseries\color{SecondaryColor} Part 2 --- the number}\\[10pt]
      {\fontsize{22}{26}\selectfont\color{PrimaryColor}
        $D=\dfrac{\log N}{\log(1/r)}$}\\[12pt]
      {\small Sierpinski $\approx 1.585$\\ Koch $\approx 1.262$}
    \end{column}
    \begin{column}{0.33\textwidth}
      \centering
      {\small\bfseries\color{SecondaryColor} Part 3 --- the world}\\[4pt]
      \includegraphics[width=0.45\textwidth]{images/britain.png}\\[4pt]
      {\small coast, lungs, maps,\\ game worlds}
    \end{column}
  \end{columns}

  \vspace{0.15cm}
  \begin{center}
    \begin{tikzpicture}
      \node[conceptbox, text width=0.85\textwidth] {
        \centering
        \textit{``Clouds are not spheres, mountains are not cones, coastlines are not
        circles\ldots''} --- Mandelbrot \\[3pt]
        \highlight{Roughness has a number, and simple rules make it.}
      };
    \end{tikzpicture}
  \end{center}
  \begin{center}{\small\color{DarkGray}\itshape Thank you --- questions welcome.}\end{center}
\end{frame}

% ---------------------------------------------------------------
% References
% ---------------------------------------------------------------
\begin{frame}{References}
  \footnotesize
  \begin{itemize}
    \item B.~Mandelbrot, ``How long is the coast of Britain? Statistical self-similarity and fractional dimension,'' \textit{Science}, 1967.
    \item B.~Mandelbrot, \textit{The Fractal Geometry of Nature}, 1982.
    \item P.~Wessel and W.~H.~F.~Smith, ``A global, self-consistent, hierarchical, high-resolution shoreline database,'' \textit{J.~Geophys.~Res.}, 1996 (GSHHG data).
    \item E.~R.~Weibel, \textit{Morphometry of the Human Lung}, Springer, 1963.
    \item M.~Häggström, ``Medical gallery of Mikael Häggström 2014,'' \textit{WikiJournal of Medicine}, 2014 (retina photograph, CC0).
    \item D.~Hilbert, ``\"Uber die stetige Abbildung einer Linie auf ein Fl\"achenst\"uck,'' \textit{Math.~Ann.}, 1891.
    \item A.~Fournier, D.~Fussell, L.~Carpenter, ``Computer rendering of stochastic models,'' \textit{Comm.~ACM}, 1982 (fractal terrain).
    \item 3Blue1Brown, ``Fractals are typically not self-similar'' (video).
  \end{itemize}
\end{frame}

\end{document}
